\chapter{Перевірка статистичних гіпотез.}

Перевірка гіпотез — це фундаментальний статистичний метод, який використовується в науці про дані для формування висновків або узагальнень про популяцію на основі вибіркових даних. Він допомагає нам визначити, чи є спостережуваний ефект у наборі даних статистично значущим, чи він міг виникнути випадково.

У науці про дані перевірка гіпотез використовується в різних сферах, зокрема в A/B-тестуванні, регресійному аналізі та оцінці моделей машинного навчання.



\section{Ключові поняття в тестуванні гіпотез}

Нульова гіпотеза ($H_0$): Нульова гіпотеза — це твердження про відсутність ефекту або різниці. Вона слугує базовим або вихідним припущенням. Наприклад, в A/B-тесті нульова гіпотеза може стверджувати, що дві групи мають однакову середню продуктивність.

Альтернативна гіпотеза ($H_1$): Альтернативна гіпотеза є протилежною нульовій гіпотезі. Вона передбачає, що існує статистично значущий ефект або різниця. Наприклад, альтернативна гіпотеза може стверджувати, що група A має вищий середній показник, ніж група B.

Тестова статистика: Тестова статистика — це стандартизоване значення, обчислене на основі вибіркових даних, яке використовується для оцінки сили доказів проти нульової гіпотези. До поширених тестових статистик належать t-оцінки та z-оцінки.

P-значення: P-значення представляє ймовірність отримання тестової статистики, що є принаймні такою ж екстремальною, як і спостережувана, припускаючи, що нульова гіпотеза є правдивою. Мале P-значення (зазвичай менше 0,05) вказує на сильні докази проти нульової гіпотези.

Рівень значущості ($\alpha$): Рівень значущості (часто позначається як $\alpha$) — це поріг, обраний аналітиком для прийняття рішення про відхилення нульової гіпотези. Зазвичай $\alpha$ встановлюється на рівні 0,05, що означає 5$\%$ ймовірність відхилення нульової гіпотези, коли вона є справжньою (помилка I типу).

Помилки типу I та типу II:

Помилка типу I: Неправильне відхилення нульової гіпотези, коли вона є істинною (хибнопозитивний результат).

Помилка типу II: Невідхилення нульової гіпотези, коли вона є хибною (хибнонегативний результат).
Потужність тесту: Потужність тесту — це ймовірність правильного відхилення нульової гіпотези, коли вона є хибною. Висока потужність означає менший ризик помилки типу II.



\textbf{Приклад}: У нас є ігровий кубик (гральна кістка).
Ми підозрюємо, що він “підкручений” — тобто шістка випадає частіше, ніж мала б.

Ми кидаємо кубик 60 разів, і отримуємо:

кількість шісток: 15

- Нульова гіпотеза $H_0$: кубик чесний
  $$
  p = \frac{1}{6} \approx 0.1667
  $$

- Альтернативна гіпотеза $H_1$: кубик “підкручений на шістку”
  $$
  p > \frac{1}{6}
  $$

(Тобто — односторонній тест.)

Випадкова величина -  Позначимо:
$$
X = \text{кількість шісток у 60 кидках}
$$
Тоді:
$$
X \sim \text{Binomial}(n = 60, p = 1/6)
$$

Очікуване значення та дисперсія тоді будуть наступні:

$$
\mathbb{E}[X] = n p = 60 \times \frac{1}{6} = 10
$$
$$
\mathrm{Var}(X) = n p (1 - p) = 60 \times \frac{1}{6} \times \frac{5}{6} \approx 8.33
$$
$$
\sigma = \sqrt{8.33} \approx 2.89
$$

Наближено нормальний тест - для (n=60) можна використати нормальне наближення:
$$
Z = \frac{X - \mathbb{E}[X]}{\sigma} = \frac{15 - 10}{2.89} \approx 1.73
$$

Рівень значущості - візьмемо рівень $\alpha = 0.05$.
Критичне значення для одностороннього тесту:
$$
Z_{кр} = 1.645
$$

Прийняття рішення - оскільки $Z = 1.73 > 1.645$, то ми відхиляємо ($H_{0}$) на рівні значущості 0.05.

Отже, є статистичні підстави вважати, що кубик дійсно “підкручений” на шістку. Це одновибірковий біноміальний тест.



\subsection{Перевірка гіпотез у машинному навчанні}

У машинному навчанні перевірка гіпотез відіграє важливу роль в оцінці моделей. Наприклад, при порівнянні двох моделей (наприклад, дерева рішень і випадкового лісу) можна використовувати перевірку гіпотез, щоб оцінити, чи є різниця в продуктивності (наприклад, точність, прецизійність) статистично значущою.

Типовим сценарієм є використання перехресного валідації для навчання моделей на різних підмножинах даних, а потім виконання статистичного тесту (наприклад, парного t-тесту) для визначення, чи одна модель постійно перевершує іншу.


\subsection{Практичний приклад: A/B-тестування з перевіркою гіпотез}

Уявіть ситуацію, коли компанія, що займається електронною комерцією, хоче перевірити, чи новий дизайн веб-сайту (B) приносить більше покупок, ніж поточний дизайн (A). Компанія проводить A/B-тестування і збирає дані про кількість покупок, зроблених користувачами, які відвідують кожну версію сайту.

Етапи:

\begin{itemize}
\item Нульова гіпотеза ($H_0$): новий дизайн не впливає на рівень покупок (тобто немає різниці між A і B).
\item Альтернативна гіпотеза ($H_1$): новий дизайн збільшує рівень покупок.
\item Рівень значущості: Встановіть $\alpha= 0.05$.
\item Збір даних: Припустимо, ви зібрали такі дані:
\item Версія A: 500 відвідувачів, 50 покупок.
\item Версія B: 600 відвідувачів, 80 покупок.
\item Проведення тесту: Використовуйте z-тест для порівняння частоти покупок між двома групами.
\item Висновок: На основі p-значення відхиліть або не відхиліть нульову гіпотезу.
\end{itemize}

Нижче ми наведемо деякі приклади статистичних перевірок.

\section{Приклади статистичних перевірок.}


\subsection{Одновибірковий t-критерій.}

Нехай у нас є дані про ріст студентів. Нашим припущенням є те, що є певний середній ріст студентів (у прикладі нижче - 183 сантиметри). Проведемо одновибірковий t-критерій (One-sample t-test), щоб перевірити цю гіпотезу.

Для цього спочатку потрібно імпортувати пакет для роботи із статистикою \verb|import scipy.stats as stats|. В метод \verb|stats.ttest_1samp| передаємо параметри: дані \verb|data| та середнє значення, яке ми припускаємо:

\begin{lstlisting}
import scipy.stats as stats 
import numpy as np
  
data = [182, 180, 190, 184]
  
t_statistic, p_value = stats.ttest_1samp(a=data, popmean=183) 
print(t_statistic , p_value)
\end{lstlisting}

Результат:

\begin{console}
0.46291004988627577 0.6749411569291093
\end{console}

Оскільки $p = 0.675 > 0.05$, ми не можемо відхилити нульову гіпотезу. Отримані дані не дають статистично значущих підстав стверджувати, що середній зріст відрізняється від 183 см.

Інакше кажучи, вибірка $182, 180, 190, 184$ узгоджується з припущенням про середнє значення 183.


\subsection{Однофакторний дисперсійний аналіз.}

Уявімо, що ми хочемо з'ясувати, чи впливає певний фактор на деякий показник. Наприклад, інтернет-магазин провів три різні рекламні кампанії, і хоче перевірити, чи вони вплинули на кількість продажів. Або ж три методики викладання впливають на успішність студентів. Постає питання: чи дійсно середні значення цих груп відрізняються, чи різниця виникла випадково?

Ідея дуже проста. Якщо середні значення всіх груп однакові, то відмінності між ними повинні бути приблизно такими ж, як випадкові коливання всередині кожної групи. Якщо ж середні істотно різняться, то різниця між групами буде значно більшою за випадковий шум. Тому ANOVA порівнює два типи дисперсії. 

Нульовою гіпотезою тут є те, що усі середні однакові. Альтернатива ($H_{1}$) - принаймні одне середнє відрізняється. 

\ding{43} Зверніть увагу: ANOVA не говорить, яка саме група відрізняється. Вона лише повідомляє, що хоча б одна середня інша.

Візьмемо дані з трьох незалежних груп, та проведемо однофакторний дисперсійний аналіз (One-way ANOVA):

\begin{lstlisting}
group1 = [2.9, 3.0, 2.5, 2.6, 3.2, 2.8]
group2 = [3.8, 2.7, 4.0, 3.9, 3.2, 3.1]
group3 = [3.1, 3.4, 3.7, 3.0, 3.1, 3.6]

f_statistic, p_value = stats.f_oneway(group1, group2, group3)

print(f"F-statistic: {f_statistic}")
print(f"P-value: {p_value}")
\end{lstlisting}

Результат:

\begin{console}
F-statistic: 4.434477379095165
P-value: 0.03068555150519455
\end{console}

Оскільки $p = 0.0307 < 0.05$ нульова гіпотеза відхиляється. Існують статистично значущі відмінності між середніми значеннями хоча б двох груп.


\subsection{Одновибірковий критерій хі-квадрат.}

Застосуємо одновибірковий критерій хі-квадрат (One-sample Chi-Squared test) до нашого прикладу із гральним кубиком. Розподіл результатів кидків грального кубика - від 1 до 6 - через масив в numpy під назвою \verb|observed|:

\begin{lstlisting}
observed = np.array([30, 10, 5, 5, 5, 5])  

avg = observed.sum()

expected = np.array([avg/6, avg/6, avg/6, avg/6, avg/6, avg/6])

chi2_stat, p_value = stats.chisquare(f_obs=observed, f_exp=expected)

print(f"Chi-Square Statistic: {chi2_stat}")
print(f"P-value: {p_value}")
\end{lstlisting}

Результат:

\begin{console}
Chi-Square Statistic: 50.0
P-value: 1.3857973367009573e-09
\end{console}

Оскільки $p << 0.05$ нульова гіпотеза впевнено відхиляється. Отриманий розподіл практично неможливо пояснити випадковими коливаннями - ігровий кубик є зваженим.


\subsection{Критерій незалежності хі-квадрат.}

Приклад даних: таблиця співвідношень зі спостережуваними частотами. Наприклад, це можуть бути дані про вподобання споживачів щодо двох брендів у різних вікових групах. Застосуємо критерій незалежності хі-квадрат (Chi-Square Test of Independence):

\begin{lstlisting}
observed = np.array([[1,2,2,5],
                     [1, 3, 2, 1]])

chi2_stat, p_value, dof, expected = stats.chi2_contingency(observed)

print(f"Chi-Square Statistic: {chi2_stat}")
print(f"P-value: {p_value}")
print(f"Degrees of Freedom: {dof}")
print(f"Expected Frequencies:\n{expected}")
\end{lstlisting}

Результат:

\begin{console}
Chi-Square Statistic: 2.412380952380953
P-value: 0.49133407367391724
Degrees of Freedom: 3
Expected Frequencies:
[[1.17647059 2.94117647 2.35294118 3.52941176]
 [0.82352941 2.05882353 1.64705882 2.47058824]]
\end{console}

Оскільки $p = 0.491 > 0.05$, ми не можемо відхилити нульову гіпотезу. Статистично значущого зв'язку між віковою групою та вибором бренду не виявлено.

Спостережувані відмінності можуть бути пояснені випадковими коливаннями вибірки.

\subsection{Коли вибрати який тест?}


\begin{table}[htbp]
\centering
\caption{Вибір статистичних критеріїв}
\label{tab:statistical_tests_grid}
\begin{tabular}{|p{3cm}|p{3.5cm}|p{4cm}|p{4cm}|}
\hline
\textbf{Тип даних / Кількість груп} & \textbf{1 група} & \textbf{2 групи} & \textbf{2 і більше груп} \\ \hline

\textbf{Категоріальні дані} & 
  Тест для частки (наближення Z-критерієм)\newline Критерій хі-квадрат & 
  Тест для частки (наближення Z-критерієм)\newline Критерій хі-квадрат & 
  Критерій хі-квадрат \\ \hline

\textbf{Неперервні дані} & 
  Z-критерій та його варіанти\newline T-критерій та його варіанти & 
  Z-критерій та його варіанти\newline T-критерій та його варіанти & 
  Дисперсійний аналіз ANOVA (F-критерій, однофакторний або двофакторний) \\ \hline

\textbf{Порушення класичних припущень} & 
  Знаковий критерій (Sign Test)\newline Критерій знакових рангів (Signed Rank Test) & 
  Критерій Вілкоксона–Манна–Уїтні\newline Парний t-критерій\newline Критерій Мак-Немара & 
  Критерій Краскела–Уолліса \\ \hline
\end{tabular}
\end{table}


\ding{46} \textbf{Питання для самоконтролю.}

\begin{enumerate}
\item Що таке статистична гіпотеза?
\item Чим нульова гіпотеза ($H_{0}$) відрізняється від альтернативної гіпотези ($H_{1}$)?
\item Яка основна мета перевірки статистичних гіпотез?
\item Що таке рівень статистичної значущості ($\alpha$)?
\item Що таке p-значення (p-value)?
\item Як приймається рішення за результатами статистичного тесту?
\item Що таке помилка першого роду (Type I Error)?
\item Що таке помилка другого роду (Type II Error)?
\item Чому перевірка статистичних гіпотез є важливою в машинному навчанні?
\item У яких задачах машинного навчання застосовують статистичні критерії?
\item Що таке A/B-тестування?
\item Які етапи проведення A/B-тесту?
\item Які показники найчастіше аналізують під час A/B-тестування?
\item Що таке одновибірковий t-критерій Стьюдента?
\item За яких умов застосовується одновибірковий t-критерій?
\item Які припущення необхідно виконати для використання t-критерію?
\item Що таке однофакторний дисперсійний аналіз (One-Way ANOVA)?
\item Для чого використовується однофакторний дисперсійний аналіз?
\item Які гіпотези перевіряє ANOVA?
\item Чому не рекомендується замість ANOVA виконувати багаторазові t-тести?
\item Що таке F-статистика в дисперсійному аналізі?
\item Що таке критерій хі-квадрат узгодженості (Goodness-of-Fit Test)?
\item У яких випадках застосовується критерій узгодженості хі-квадрат?
\item Що таке очікувані та спостережувані частоти?
\item Що перевіряє критерій незалежності хі-квадрат?
\item Які дані необхідні для застосування критерію незалежності хі-квадрат?
\item Що таке таблиця спряженості (Contingency Table)?
\item Які припущення необхідно виконати для застосування критеріїв хі-квадрат?
\item Як інтерпретувати результати перевірки статистичної гіпотези?
\item Які обмеження слід враховувати під час використання статистичних критеріїв у задачах інтелектуального аналізу даних?
\end{enumerate}